\begin{frame}{채점 기준: ``채웠는가''가 아니라 ``구체적인가''}
  \small
  \begin{tabular}{@{}llp{7.8cm}@{}}
    \toprule
    품질 & 특징 & 예시 \\
    \midrule
    1점 --- 감상평 & 구체적 지적이 없음, 느낌만 &
      ``발표 잘 들었습니다. 결과가 인상적이네요.'' \\
    2점 --- 일반 지적 & 문제를 제기하지만 \textbf{근거(개념)가
      약함} --- 발표팀이 ``네''로 넘길 수 있음 &
      ``모델 선정 이유가 부족해 보입니다. 데이터 특성을 더
      봐주세요.'' \\
    3점 --- 반박 가능 & \textbf{장·절을 명시}하고, 답이 나쁘면
      결과가 흔들리는 질문 & ``불균형(99.8:0.2)에 정확도만
      비교했는데(Ch02.3), Precision·Recall을 안 주면 '둘 다
      항상 정상으로 찍는' 모델일 수 있는데 어떻게 구분
      하셨나요?'' \\
    \bottomrule
  \end{tabular}
\end{frame}
